\section{Supplementary analyses and execution record}

\subsection{Scope and provenance}

This supplement accompanies the revised audit. It does not constitute a public preregistration. The internal design and worklog were maintained locally; no independent public timestamp establishes when each choice was fixed. Primary slice-level intervals retain the original sealed predictions and fixed split; all revision-stage analyses below are exploratory.

\subsection{The thirteen libraries and the audited set}

The release carries thirteen library directories under one publishing organisation and the audited
article documents six datasets. Table~\ref{tab:libraries} lists the nine libraries this audit
touches and what each was used for. Three of the five main-sample libraries --- CTIS, CNPM and bel canto --- are
among the six datasets the audited article evaluates. Music genre and pianos belong to the same
published collection but are not covered by that article, so contrasts in them audit the same
kind of split under a framework that does not report on them. The three negative controls each
contribute one slice per recording, with GZ\_IsoTech carrying 46 recordings of two to four
segments. Guzheng\_Tech99 was excluded because every recording contributes one slice and its
per-note annotations have the same per-recording dominant label in all 99 recordings, so no
recording-level class layer exists. The four remaining libraries of the release were not part of
the audited set; we make no claim about them.

\begin{table}[ht]\centering\footnotesize\setlength{\tabcolsep}{4pt}
\caption{Libraries touched by this audit. Slices, recordings and classes are counted in the released MEL
manifest after canonicalisation of duplicate recordings. The frozen CTIS panel keeps 201 of the 219
classes and 4{,}887 of the 4{,}954 recordings after the eligibility rule and the placeholder removal;
the other four main-sample libraries keep all their classes. Guzheng\_Tech99 releases 99 MEL images
without a per-recording class directory, so its class count is not reported.}
\label{tab:libraries}
\begin{tabular}{@{}lrrrl@{}}\toprule
Library & Slices & Recordings & Classes & Role in this audit\\\midrule
CTIS & 43{,}054 & 4{,}954 & 219 & primary; A0, A1, A2, C1\\
music genre & 36{,}375 & 1{,}713 & 16 & primary; A0, A1, A2\\
pianos & 18{,}352 & 668 & 8 & primary; A0, A1, A2\\
CNPM & 2{,}724 & 287 & 5 & primary; A0, A1, A2\\
bel canto & 9{,}910 & 203 & 4 & primary; A0, A1, A2\\
chest\_falsetto & 1{,}280 & 1{,}280 & 4 & negative control; A0, A2\\
erhu\_playing\_tech & 1{,}253 & 1{,}253 & 11 & negative control; A0, A2\\
GZ\_IsoTech & 2{,}899 & 2{,}824 & 8 & negative control; A0, A2\\
Guzheng\_Tech99 & 99 & 99 & --- & excluded; no recording-level class layer\\
\bottomrule\end{tabular}\end{table}

\subsection{Execution changes}

The list below records every difference between the internal plan and the executed
runs, and adds a column the earlier version did not carry: what was known when the decision was
taken. The column matters because a change made with the relevant count in hand is a different
object from a change made after seeing a result. Nothing in the list is a claim that the
sequence was planned; it records what was available, and the supplementary analyses that follow
are marked exploratory where they were added at revision.

\begin{enumerate}\itemsep2pt
\item \textbf{CTIS classes:} 219 $\rightarrow$ 201. Known at the decision: per-class recording counts and the no-audio placeholder flags. Reason: the five-recording eligibility rule, then placeholder removal. Effect: 8.2\% fewer classes; 545 slices (1.3\%) drop.
\item \textbf{Injection volume:} uniform 50\%/30\% of $S_c$ $\rightarrow$ all $Q^{\text{pool}}$. Known at the decision: the capacity calculation on the 201 classes: 10 can supply 50\%, 38 can supply 30\%. Reason: a common fraction is infeasible for most classes. Effect: $\alpha_c$ varies by class, median 0.17; no fixed common exposure.
\item \textbf{Exposure agreement:} enforce within 2\% of a 400-draw mean $\rightarrow$ report descriptively. Known at the decision: 400-draw planning means and per-class capacity. Reason: the C-side recording count is the binding constraint. Effect: $\alpha_c$ dispersion is reported, not tested.
\item \textbf{Negative controls:} one seed $\rightarrow$ three seeds. Known at the decision: the initial single-seed contrasts. Reason: single-seed dispersion is too wide to read. Effect: exploratory; five of six intervals include zero.
\item \textbf{Slice budget:} $m_c=0.60S_c$ $\rightarrow$ capped by C capacity. Known at the decision: the classwise feasibility artifact. Reason: C cannot supply the planned count. Effect: a slice-count constraint, not early stopping.
\item \textbf{C1 control:} three arms $\rightarrow$ three arms plus C1. Known at the decision: the main-arm results. Reason: separate content from recording identity without test recordings. Effect: descriptive only; validation not interpretable.
\item \textbf{Cross-instrument head:} L-inst $\rightarrow$ L-inst plus adversarial head, weight $\le 0.02$. Known at the decision: the instrument-probe pilot. Reason: instrument identity survives in the representation. Effect: the adversarial arm is exploratory.
\item \textbf{Probe split:} 80/20 of all recordings $\rightarrow$ 80/20 of training recordings. Known at the decision: withheld instruments are absent from the training label space. Reason: a closed-set probe cannot identify unseen classes. Effect: the probe measures identity, not recognition.
\item \textbf{Analyses not run:} selection regret; CQT/chroma $\rightarrow$ not run. Known at the decision: per-epoch sealed-panel scores were not stored. Reason: no artifact to analyse. Effect: stated as not run.
\item \textbf{Recording-level estimator:} slice-weighted point $\rightarrow$ one vote per recording. Known at the decision: human review found that the point and the intervals used different weights. Reason: align the point estimator with the intervals. Effect: ten points change, all inside the retained intervals.
\item \textbf{Revision-stage additions:} none $\rightarrow$ class-stratified and joint bootstrap, self-test tables, version comparisons, partial compute ledger. Known at the decision: reviewer comments. Reason: transparency. Effect: exploratory.
\end{enumerate}

\subsection{Realised injection and training steps}

Table~\ref{tab:realised-injection} counts pool slices actually present in A1 training. The exposed-test numerator counts canonical recording identifiers appearing in both A1's injected pool slices and the test panel; the denominator counts unique test recording identifiers. Injected recordings also include pool recordings absent from the test panel. Steps are calculated per model as ten times the integer number of full 64-slice batches, with the last partial batch dropped; both backbones use the same counts. Counts are derived from the frozen split CSVs, not independent evidence of the timing of the split.

\begin{table}[htbp]\centering\small
\caption{Actual A1 injection and per-model optimiser steps. Test exposure is exposed test recordings / all test recordings.}
\label{tab:realised-injection}
\begin{tabular}{lrrrrr}\toprule
Library & Injected slices & Injected rec. & Exposed test rec. & A0 steps & A1 steps \\\midrule
CTIS & 4{,}178 & 753 & 578/963 & 3{,}980 & 3{,}980 \\
music genre & 3{,}613 & 340 & 335/337 & 3{,}410 & 3{,}410 \\
pianos & 1{,}862 & 134 & 133/135 & 1{,}720 & 1{,}720 \\
CNPM & 268 & 50 & 47/54 & 250 & 250 \\
bel canto & 977 & 40 & 40/40 & 920 & 920 \\
\bottomrule\end{tabular}\end{table}

\subsection{Convergence at the end of training}

Per-epoch training logs were retained for every run, so the arms can be compared where they
stop. On CTIS the final-epoch training loss has median 0.9016 in A0, 0.9027 in A1 and 0.9026 in
C1 under AlexNet, and 0.8712, 0.8715 and 0.8714 under Swin-T; final-epoch training accuracy is at
least 0.9997 in every arm. A2, which has 25{,}263 training slices and 3{,}940 steps against
3{,}980, ends at 0.9034 and 0.8718. The arms therefore end at the same point of the objective,
consistent with equal epochs and, for A0, A1 and C1, equal step counts. The logs hold
training-set and C-side validation metrics only: per-epoch sealed-panel scores were not stored,
so no selection-regret curve is reported.

\subsection{CTIS class exclusions}

The source manifest contains 219 classes. Thirteen have fewer than five source recordings:
C0183 (4), C0200 (2), C0201 (1), C0257 (3), D0049 (4), D0050 (4), D0103 (4),
D0104 (4), D0246 (4), D0279 (4), D0284 (2), D0287 (4), and D0299 (4).
Five further classes, D0051, D0058, D0060, D0061 and D0062, each have five source
recordings but fall to four after removal of one recording whose identifier contains the
Chinese no-audio marker. These exclusions leave 201 classes. The five-recording threshold
is a study-specific eligibility rule, not a requirement of every recording-level split.

\subsection{Threshold sensitivity and the profile of excluded classes}

Because the five-recording rule is a study-specific choice, we report what other thresholds
would do. Every row below applies the threshold and then removes the five placeholder classes,
so each row is comparable with the frozen 201-class set. Thresholds of 2, 3, 5 and 10 source
recordings keep 213, 211, 201 and 159 classes, covering 42{,}882, 42{,}834, 42{,}509 and
41{,}035 of the 43{,}054 released slices and 1{,}272, 1{,}270, 1{,}241 and 1{,}035 Q
recordings. The frozen set sits where a stricter threshold starts to remove classes in bulk:
from 5 to 10 the class count falls by 42 while the slice count falls by 1{,}474, so the
threshold trades many small classes for few slices.

The classes removed by the rule are smaller than the retained classes in absolute size but not
in density: the thirteen classes with fewer than five recordings have median 4 source
recordings against 23, median 20 slices against 197, and 8.25 slices per recording against 6.90.
The rule removes 393 slices, 0.9\% of the library; the five placeholder classes remove a
further 152 slices, 0.35\%. Together the eighteen excluded classes account for 545 slices,
1.3\% of the library, so the exclusion does not carve out a systematically harder stratum.
Fourteen retained classes have four to six source recordings; their recording-level intervals
are the widest in the primary results and are reported with the rest, without a separate
stability claim.

\subsection{Cross-library duplication check}

The same slice could in principle be registered under two libraries. We compared the content
hash of every MEL slice across the five libraries: 110{,}415 slices, 110{,}280 unique hashes,
and no hash appears in two libraries. The 135 repeated instances are within-library, consistent
with the 77 identical-image groups reported for CTIS. The check is exact byte identity only:
re-encoded or re-cropped copies of the same audio would not be caught, and a recording-level
cross-library comparison is not possible because recording identifiers are defined independently
per library. The check is reproduced by \texttt{code/audit\_cross\_library\_duplicates.py},
which writes \texttt{data/results/cross\_library\_duplicates.json}.

\subsection{A2 self-test and contextual published scores}

Self-test scores are means and sample standard deviations over three seeds. CNPM Swin-T has a self-test standard deviation of 0.0366; the available results do not establish its cause. Their difference from A0 mixes different evaluation panels and is not a paired pool-injection contrast. Published scores are contextual comparisons only: the original split files, training realisation and retained label universe differ. A dash means no verified matching published cell. Published CTIS values are from Table 3 (printed p.33), CNPM from Table 7 and singing style from Table 8 (printed p.35) of DOI 10.5334/tismir.194. Those cells all report mel-input F1 under the article's own split, optimizer and input pipeline, which differ from ours on every axis (the article uses ImageNet V1 weights, SGD at 0.001, 40 epochs, a batch of 4 and weighted cross-entropy), so the column locates our numbers in the published landscape and is not a reproduction check. The CTIS table also lists 0.936 for a ViT-L/32 mel model, so the same figure attaches to two rows there; we quote the AlexNet row.

\begin{table}[ht]\centering\small

\caption{Self-test scores and contextual published MEL F1. The difference column is in percentage points, not a causal decomposition.}

\begin{tabular}{llrrrr}\toprule

Library & Backbone & Self-test & s.d. & Self$-$A0 (cross-panel) & Published \\\midrule

CTIS & AlexNet & 0.9645 & 0.0060 & +5.67 & 0.936 \\

CTIS & Swin-T & 0.9851 & 0.0016 & +3.15 & 0.956 \\

music genre & AlexNet & 0.8988 & 0.0023 & +20.99 & --- \\

music genre & Swin-T & 0.9515 & 0.0039 & +21.69 & --- \\

pianos & AlexNet & 0.9803 & 0.0057 & +15.66 & --- \\

pianos & Swin-T & 0.9916 & 0.0039 & +12.17 & --- \\

CNPM & AlexNet & 0.7708 & 0.0175 & +36.43 & 0.742 \\

CNPM & Swin-T & 0.8349 & 0.0366 & +43.47 & --- \\

bel canto & AlexNet & 0.9314 & 0.0048 & +11.14 & 0.919 \\

bel canto & Swin-T & 0.9746 & 0.0009 & +16.55 & 0.906 \\

\bottomrule\end{tabular}\end{table}

\subsection{Instrument-identity probe}

These are means of the fold-level one-vs-rest macro-AUC and top-1 instrument accuracy, then averaged over two probe-fit seeds. The probe's 80/20 split of the training recordings uses seed 0; the two seeds below are seeds of the probe fit, not of the classifier training. They are not technique classification scores.

\begin{table}[ht]\centering\small\caption{Training-side instrument probe; two probe-fit seeds, three folds, closed class set.}

\begin{tabular}{llrr}\toprule Arm & Backbone & Macro-AUC & Top-1 \\\midrule

R-slice & AlexNet & 0.9324 & 0.5550 \\

R-slice & Swin-T & 0.9394 & 0.6326 \\

L-inst & AlexNet & 0.9612 & 0.6682 \\

L-inst & Swin-T & 0.9697 & 0.7320 \\

L-inst+Adv & AlexNet & 0.9261 & 0.5952 \\

L-inst+Adv & Swin-T & 0.9366 & 0.6693 \\

\bottomrule\end{tabular}\end{table}

\subsection{Fold composition}
\label{sec:fold-composition}

The three arms that report fold-level results use different fold definitions, because the fold has to
be the unit the arm is protecting: the slice-level arm A0 uses random slice folds, the exposure and
reconstruction arms A1 and A2 share instrument-grouped folds, and the recording-grouped arm A1r uses
recording folds. The numbers below are the realised composition of each fold, recomputed from the
frozen split files and the stored fold results.

\begin{table}[ht]\centering\small\caption{Realised fold composition of the cross-instrument experiment (CTIS). Test slices are the evaluated panel of each fold.}
\label{tab:fold-composition}
\begin{tabular}{llrrrr}\toprule Folds & Fold & Train rec. & Test rec. & Test slices & Classes \\ \midrule
A0, slice & 1 & 1{,}151 & 1{,}019 & 2{,}524 & 43 \\
A0, slice & 2 & 1{,}146 & 1{,}014 & 2{,}508 & 43 \\
A0, slice & 3 & 1{,}164 & 997 & 2{,}496 & 43 \\
A1/A2, instrument & 1 & 829 & 467 & 2{,}509 & 42 \\
A1/A2, instrument & 2 & 818 & 480 & 2{,}509 & 43 \\
A1/A2, instrument & 3 & 853 & 441 & 2{,}510 & 43 \\
A1r, recording & 1 & 835 & 461 & 2{,}509 & 43 \\
A1r, recording & 2 & 834 & 462 & 2{,}509 & 43 \\
A1r, recording & 3 & 831 & 465 & 2{,}510 & 43 \\
\bottomrule\end{tabular}\end{table}

Two features of the table matter for reading Section~6.4. First, the slice-level folds of A0 place
slices of the same recording on both sides, which is exactly the leak the arm is meant to measure;
the instrument-grouped and recording-grouped folds share no recording. Second, the first
instrument-grouped fold covers 42 of
the 43 classes on its test side, because the natural-harmonic technique has no test recording there;
the pooled technique score averages over the classes present, so that fold contributes 42 classes.
The 135 instruments each fall in exactly one instrument-grouped fold.

\subsection{Recording-level estimator correction}

The old point estimator weighted the averaged prediction of each recording by its slice count. The original recording-cluster intervals instead counted each recording once. The revision makes the point estimator consistent with those intervals; slice-level results are unchanged. The following are A1--A0 point differences in percentage points. Original intervals are retained, because their bootstrap already used equal recording votes. All ten corrected point estimates lie within their respective retained intervals. In bel canto the corrected values exceed the old values by 4.43 and 7.54 points for AlexNet and Swin-T, respectively. Different weighting defines a different estimand; the size of the change does not by itself identify its cause. One library is also constrained by a ceiling: in bel canto the exposure arm reaches 1.00 at the recording level for most seeds against 0.84--0.93 in the control, so its recording-level contrast and interval are limited by the metric's upper bound as well as by its 40 test recordings.

\begin{table}[ht]\centering\small\caption{Corrected recording-level points; old slice-weighted points shown only to document the change.}

\begin{tabular}{llrr}\toprule Library & Backbone & Old & Corrected \\ \midrule

CTIS & AlexNet & +1.67 & +1.37 \\

CTIS & Swin-T & +1.64 & +1.61 \\

CNPM & AlexNet & +26.45 & +28.44 \\

CNPM & Swin-T & +29.91 & +30.75 \\

bel canto & AlexNet & +10.18 & +14.61 \\

bel canto & Swin-T & +15.20 & +22.74 \\

music genre & AlexNet & +16.07 & +15.13 \\

music genre & Swin-T & +16.63 & +15.98 \\

pianos & AlexNet & +7.63 & +7.53 \\

pianos & Swin-T & +4.85 & +4.86 \\

\bottomrule\end{tabular}\end{table}

\subsection{CTIS uncertainty sensitivity}

These exploratory intervals use an independently implemented bootstrap (2,000 resamples; fixed split). The point estimate is always computed on the observed data, not the bootstrap mean. Class stratification resamples recording clusters within true class. Joint resampling also draws the three trained seeds with replacement, preserving arm pairing. Three observed seeds cannot identify uncertainty for a new training population. The original primary intervals remain in the manuscript; these sensitivity results are not selected to replace them. The independent reimplementation reproduces the CTIS point estimates exactly and the cluster-only intervals to within a few hundredths: 2.02 with bounds 0.79 to 3.28 here against 0.83 to 3.35 in the sealed pipeline, and 1.70 with bounds 1.27 to 2.81 against 1.28 to 2.84 there. The manuscript reports the sealed-pipeline interval; both AlexNet upper bounds exceed the three-point threshold. The small difference is consistent with the different resampling stream, not a change of estimate.

\begin{table}[ht]\centering\scriptsize\caption{CTIS sensitivity: percentage-point effects and 95\% percentile intervals.}

\begin{tabular}{lllrrr}\toprule Metric & Backbone & Scheme & Point & Lower & Upper \\ \midrule

slice & AlexNet & cluster, fixed seeds & +2.02 & +0.79 & +3.28 \\

slice & AlexNet & class, fixed seeds & +2.02 & +1.02 & +2.84 \\

slice & AlexNet & class + seeds & +2.02 & +0.57 & +3.30 \\

slice & AlexNet & cluster + seeds & +2.02 & +0.35 & +3.74 \\

recording & AlexNet & cluster, fixed seeds & +1.37 & +0.28 & +2.75 \\

recording & AlexNet & class, fixed seeds & +1.37 & +0.43 & +2.38 \\

recording & AlexNet & class + seeds & +1.37 & +0.01 & +2.73 \\

recording & AlexNet & cluster + seeds & +1.37 & -0.03 & +2.99 \\

slice & Swin-T & cluster, fixed seeds & +1.70 & +1.27 & +2.81 \\

slice & Swin-T & class, fixed seeds & +1.70 & +1.33 & +2.21 \\

slice & Swin-T & class + seeds & +1.70 & +0.83 & +2.59 \\

slice & Swin-T & cluster + seeds & +1.70 & +0.91 & +3.14 \\

recording & Swin-T & cluster, fixed seeds & +1.61 & +0.86 & +2.63 \\

recording & Swin-T & class, fixed seeds & +1.61 & +1.13 & +2.11 \\

recording & Swin-T & class + seeds & +1.61 & +0.61 & +2.84 \\

recording & Swin-T & cluster + seeds & +1.61 & +0.50 & +3.13 \\

\bottomrule\end{tabular}\end{table}

The AlexNet recording-level cluster-plus-seed interval includes zero (approximately -0.03 to +2.99). Both primary slice-level effects remain positive in these sensitivity analyses, but some upper bounds exceed the three-point threshold, including Swin-T under unstratified joint resampling. Thus neither a universal non-zero recording-level effect nor a universal below-SESOI conclusion is warranted.

\subsection{Macro-F1 class-set convention}

Macro-F1 needs a decision about classes that a resampling replicate does not contain. The primary
analysis recomputes the class set inside each replicate and drops a class absent from that
replicate; a fixed-class-set variant instead fixes the set once from the observed panel and scores
an absent class as zero. The point estimates are identical under both conventions, and the intervals
differ by at most 0.4 points (Table~\ref{tab:classset}): A1--A0 under AlexNet is +2.02 with bounds +0.72 to +3.09 under a fixed
class set, against +0.83 to +3.35 as reported; A1--C1 under Swin-T is +1.07 with bounds +0.20 to
+1.71 under a fixed class set, against +0.23 to +1.96 as reported. Under the recomputed set a
replicate loses on average about 7 of the 201 CTIS classes (at most 18), so the two conventions answer
slightly different questions; the manuscript and the table above use the recomputed set throughout.
No conclusion in the paper changes between them.

\begin{table}[t]
\centering
\caption{CTIS effects under the recomputed and the fixed class set (percentage points). Both
conventions score the same sealed panel and use the same 2{,}000 recording-cluster replicates;
the point estimate is identical because it is computed on the observed panel, so only the tails of
the percentile interval move. The full-class-set convention gives the same numbers as the fixed
one and is omitted.}
\label{tab:classset}
\small
\begin{tabular}{llrrr}
\toprule
Backbone & Contrast & Class set & Point & 95\% interval \\
\midrule
AlexNet & A1$-$A0 & recomputed & +2.02 & $[+0.81,+3.36]$ \\
AlexNet & A1$-$A0 & fixed      & +2.02 & $[+0.72,+3.09]$ \\
AlexNet & A1$-$C1 & recomputed & +1.96 & $[+0.79,+2.91]$ \\
AlexNet & A1$-$C1 & fixed      & +1.96 & $[+0.87,+2.65]$ \\
Swin-T  & A1$-$A0 & recomputed & +1.70 & $[+1.28,+2.81]$ \\
Swin-T  & A1$-$A0 & fixed      & +1.70 & $[+0.93,+2.36]$ \\
Swin-T  & A1$-$C1 & recomputed & +1.07 & $[+0.21,+1.93]$ \\
Swin-T  & A1$-$C1 & fixed      & +1.07 & $[+0.20,+1.71]$ \\
\bottomrule
\end{tabular}
\end{table}

\subsection{Finite resampling tail convention}

The legacy script clipped a zero-tail two-sided estimate at 1/2,000. Both primary CTIS tails had no non-positive resampled differences under that original stream. Applying the add-one convention gives raw 2/2,001 and Holm-adjusted 4/2,001 for each of the two family tests (approximately 0.0010 and 0.0020). This documents a finite-resolution convention, not an exact null probability.

\subsection{Partial compute ledger}

Integer-minute driver logs represent batch wall-clock time while training on one GPU, not metered GPU occupancy. Known batch entries are approximately 6.25 h for A0/A1 (60 models), 4.15 h for A2 (30 models), 2.99 h for the four-arm cross-instrument set (48 fold-models), 1.17 h for C1 (six models) and 0.55 h for a recording-grouped/probe driver. The last entry mixes training and probe work; some intermediate-arm work is already represented in the fold ledger, so entries must not be blindly added. Negative-control and abandoned-run logs are incomplete. No exact study-wide GPU-hour total is supported.

\subsection{Reproducible commands}

Analysis uses Python 3.10.22, PyTorch 2.11.0+cu128 and one RTX 5060 Ti 16 GB. Predictions store probabilities, labels and recording identifiers. Original training is not re-run for this revision. Run the following from the project root:

\begin{verbatim}

python code/human_review_selftest.py
python code/audit_ctis_classes.py
python code/audit_compute_ledger.py
python code/audit_ctis_bootstrap.py
python code/audit_review_exposure_20261108.py
python code/round2_a1_minus_c1.py
python code/round2_macro_denominator.py
python code/round2_data_accounting.py
python code/verify_paper_numbers.py

\end{verbatim}

\subsection{Availability and remaining limits}

The underlying recordings/images are not redistributed. This supplement supplies analysis summaries for review, not the complete prediction archive. Manifests, splits, predictions and analysis code remain available from the corresponding author on request, with public archival deposit planned upon acceptance. Only one split realisation and three primary training seeds are available; no bootstrap creates evidence about new split realisations. Compute timing logs are incomplete and rounded; no exact full-study GPU-hours are claimed.
